%%%%%%%%%%%%%%%%%%%%%%%%%%%%%%%%%%%%%%%%%%%%%%%%%%%%%%%%%
%
%           GUIA DE ESTUDO: TEOREMAS LIMITE
%           Baseado na 4ª Lista de CE308
%
%%%%%%%%%%%%%%%%%%%%%%%%%%%%%%%%%%%%%%%%%%%%%%%%%%%%%%%%%

% --- Configurações do Documento ---
\documentclass[a4paper, 10pt]{article} % <-- MUDANÇA 1: Fonte para 10pt (o padrão é 10, 11 ou 12)

% --- Pacotes Essenciais ---
\usepackage[utf8]{inputenc}
\usepackage[T1]{fontenc}
\usepackage[portuguese]{babel}
\usepackage[a4paper, margin=1.5cm]{geometry} % <-- MUDANÇA 2: Margens BEM pequenas (1.5cm)
\usepackage{amsmath, amssymb, amsthm}
\usepackage{graphicx}
\usepackage{hyperref}
\usepackage{xcolor}

% --- Otimização Agressiva de Espaço ---

% \usepackage{parskip}                % Mantém desabilitado (economiza espaço)
\setlength{\parindent}{15pt}         % Mantém a indentação do parágrafo

% MUDANÇA 3: Diminui o espaço entre linhas
\usepackage{setspace}
\setstretch{0.9} % Reduz o espaçamento entre linhas em 10%. (Tente 0.85 se ainda for muito)

% MUDANÇA 4: Diminui o espaço ANTES e DEPOIS dos títulos de seção
\usepackage{titlesec}
\titlespacing*{\section}{0pt}{1.5ex}{1ex} % (espaço antes, espaço depois)
\titlespacing*{\subsection}{0pt}{1.5ex}{1ex}
\titlespacing*{\subsubsection}{0pt}{1.5ex}{1ex}

% MUDANÇA 5: Diminui o espaço ANTES e DEPOIS de equações display ($$)
\setlength{\abovedisplayskip}{4pt}      % Padrão é ~10pt
\setlength{\belowdisplayskip}{4pt}      % Padrão é ~10pt
\setlength{\abovedisplayshortskip}{0pt}
\setlength{\belowdisplayshortskip}{0pt}


% --- Definição de Ambientes (Teoremas, Dicas, etc.) ---
\theoremstyle{definition}
\newtheorem{definition}{Definição}[section] % Reinicia a contagem por seção

\theoremstyle{plain}
\newtheorem{theorem}{Teorema}[section]

\theoremstyle{remark}
\newtheorem{example}{Exemplo}[section]
\newtheorem*{tip}{Dica de Estudo} % Dica sem numeração

% --- Comandos Personalizados (para atalhos) ---
\newcommand{\R}{\mathbb{R}}
\newcommand{\N}{\mathbb{N}}
\newcommand{\E}{\mathbb{E}}
\newcommand{\Var}{\text{Var}}
\newcommand{\Cov}{\text{Cov}}
\newcommand{\Prob}{\mathbb{P}}
\newcommand{\pto}{\xrightarrow{P}}      % Convergência em Probabilidade
\newcommand{\dto}{\xrightarrow{D}}      % Convergência em Distribuição
\newcommand{\qcto}{\xrightarrow{q.c.}}  % Convergência Quase Certa
\newcommand{\Xn}{\bar{X}_n}             % Média amostral

% --- Informações do Título ---
\title{Guia de Estudo Completo \\ Convergência e Teoremas Limite \\ \large (Baseado na Lista 4 de CE308)}
\author{Miqueias}
\date{\today}

%%%%%%%%%%%%%%%%%%%%%%%%%%%%%%%%%%%%%%%%%%%%%%%%%%%%%%%%%
% --- INÍCIO DO DOCUMENTO ---
%%%%%%%%%%%%%%%%%%%%%%%%%%%%%%%%%%%%%%%%%%%%%%%%%%%%%%%%%
\begin{document}

\maketitle
\tableofcontents
\newpage

\section{Os Modos de Convergência}
\label{sec:modos}
O objetivo desta seção é entender as diferentes formas como uma sequência de variáveis aleatórias ($X_n$) pode "convergir" para outra variável ($X$).

% ---------------------
\subsection{Convergência em Distribuição (\texorpdfstring{$X_n \dto X$}{X\_n -> D X})}
\label{sec:conv_dist}

\begin{definition}
Dizemos que $X_n$ converge em distribuição para $X$ se a Função de Distribuição Acumulada (FDA) de $X_n$ converge para a FDA de $X$, em todos os pontos onde a FDA de $X$ é contínua.
$$ \lim_{n \to \infty} F_n(x) = F(x) \quad (\forall x \text{ onde } F \text{ é contínua}) $$
\end{definition}

\subsubsection{Como Provar?}
Existem dois métodos principais:

\begin{enumerate}
    \item \textbf{Método da FDA (Direto):}
        \begin{enumerate}
            \item Encontre a FDA $F_n(x) = \Prob(X_n \le x)$.
            \item Calcule o limite $F(x) = \lim_{n \to \infty} F_n(x)$.
            \item Identifique a distribuição $X$ que possui $F(x)$ como sua FDA.
        \end{enumerate}

    \item \textbf{Método das Funções Geradoras (Teorema da Continuidade de Lévy):}
        Se a Função Geradora de Momentos (MGF) ou a Função Característica (FC) de $X_n$ converge para a MGF/FC de $X$, então $X_n \dto X$.
        $$ \lim_{n \to \infty} M_{X_n}(t) = M_X(t) \implies X_n \dto X $$
        (Isso é especialmente útil para somas de V.A.s, como no TCL).
\end{enumerate}

\begin{tip}
    A convergência em distribuição é o tipo de convergência \textbf{mais fraco}. Ela diz que as \textit{probabilidades} se parecem, mas não diz nada sobre os \textit{valores} de $X_n$ e $X$ em si.
\end{tip}

\begin{example}[Exercício 3]
    Dado $X_n \sim U(0, 1 + 1/n)$.
    \begin{enumerate}
        \item \textbf{FDA $F_n(x)$:} $F_n(x) = \frac{x - 0}{1 + 1/n - 0} = \frac{x}{1 + 1/n}$, para $x \in [0, 1+1/n]$.
        \item \textbf{Limite:} $\lim_{n \to \infty} F_n(x) = \lim_{n \to \infty} \frac{x}{1 + 1/n} = \frac{x}{1+0} = x$.
        \item \textbf{Identificação:} A FDA $F(x) = x$ (para $x \in [0,1]$) é a FDA de uma $U(0, 1)$.
    \end{enumerate}
    Portanto, $X_n \dto U(0, 1)$.
\end{example}

\begin{example}[Exercício 7]
    Dado $X_n \sim Exp(\frac{2n}{n+1})$. O parâmetro $\lambda_n = \frac{2n}{n+1} \to 2$.
    \begin{enumerate}
        \item \textbf{FDA $F_n(x)$:} $F_n(x) = 1 - e^{-\lambda_n x} = 1 - e^{-(\frac{2n}{n+1})x}$.
        \item \textbf{Limite:} $\lim_{n \to \infty} F_n(x) = 1 - e^{-(\lim \lambda_n)x} = 1 - e^{-2x}$.
        \item \textbf{Identificação:} $F(x) = 1 - e^{-2x}$ é a FDA de $X \sim Exp(2)$.
    \end{enumerate}
    Portanto, $X_n \dto X \sim Exp(2)$.
\end{example}

% ---------------------
\subsection{Convergência em Probabilidade (\texorpdfstring{$X_n \pto X$}{X\_n -> P X})}
\label{sec:conv_prob}

\begin{definition}
Dizemos que $X_n$ converge em probabilidade para $X$ se, para qualquer "margem de erro" $\epsilon > 0$, a probabilidade de $X_n$ estar longe de $X$ vai para zero.
$$ \lim_{n \to \infty} \Prob(|X_n - X| > \epsilon) = 0 \quad (\forall \epsilon > 0) $$
\end{definition}

\subsubsection{Como Provar?}
\begin{enumerate}
    \item \textbf{Método da Definição (Direto):}
        \begin{enumerate}
            \item Isole o evento: $|X_n - X| > \epsilon$.
            \item Calcule a probabilidade $p_n = \Prob(|X_n - X| > \epsilon)$.
            \item Mostre que $\lim_{n \to \infty} p_n = 0$.
        \end{enumerate}
    \item \textbf{Método de Chebyshev / $L_2$:}
        Se $\lim_{n \to \infty} \E[(X_n - X)^2] = 0$ (convergência em $L_2$), então $X_n \pto X$.
        (Isso vem da Desigualdade de Chebyshev: $\Prob(|X_n - X| > \epsilon) \le \frac{\E[(X_n - X)^2]}{\epsilon^2}$).
\end{enumerate}

\begin{tip}
    Se $X_n$ converge em probabilidade para uma \textbf{constante} $c$ (ou seja, $X_n \pto c$), então $X_n$ também converge em \textit{distribuição} para $c$ (uma V.A. degenerada).
\end{tip}

\begin{example}[Exercício 1]
    Provar $Y_n = X_n / \log n \pto 0$, onde $X_n \sim Exp(1)$.
    \begin{enumerate}
        \item \textbf{Evento:} Queremos $\lim \Prob(|Y_n - 0| > \epsilon) = 0$.
        $|X_n / \log n| > \epsilon \implies X_n > \epsilon \log n$ (pois $X_n > 0$ e $n$ grande).
        \item \textbf{Probabilidade:} $\Prob(Y_n > \epsilon) = \Prob(X_n > \epsilon \log n)$.
        Para $X_n \sim Exp(1)$, a função de sobrevivência é $\Prob(X > x) = e^{-x}$.
        $$ p_n = \Prob(X_n > \epsilon \log n) = e^{-\epsilon \log n} = e^{\log(n^{-\epsilon})} = n^{-\epsilon} = \frac{1}{n^\epsilon} $$
        \item \textbf{Limite:} $\lim_{n \to \infty} p_n = \lim_{n \to \infty} \frac{1}{n^\epsilon}$. Como $\epsilon > 0$, $n^\epsilon \to \infty$, e o limite é 0.
    \end{enumerate}
    Portanto, $Y_n \pto 0$.
\end{example}

% ---------------------
\subsection{Convergência Quase Certa (\texorpdfstring{$X_n \qcto X$}{X\_n -> q.c. X})}
\label{sec:conv_qc}

\begin{definition}
(A mais forte). Dizemos que $X_n$ converge quase certamente para $X$ se o conjunto de "mundos" (outcomes $\omega$) onde $X_n(\omega)$ não converge para $X(\omega)$ tem probabilidade zero.
$$ \Prob\left( \lim_{n \to \infty} X_n = X \right) = 1 $$
\end{definition}

\subsubsection{Como Provar?}
Provar isso diretamente é difícil. Geralmente usamos:
\begin{enumerate}
    \item \textbf{Lei Forte dos Grandes Números (SLLN):} (Veja Seção \ref{sec:lln}).
    \item \textbf{Lemas de Borel-Cantelli:}
        Seja $A_n$ o evento "ruim" (ex: $|X_n - X| > \epsilon$).
        \begin{itemize}
            \item \textbf{1º Lema:} Se $\sum_{n=1}^\infty \Prob(A_n) < \infty$, então $\Prob(A_n \text{ i.o.}) = 0$. (Os eventos ruins acontecem um número finito de vezes, então a convergência ocorre).
            \item \textbf{2º Lema:} Se $A_n$ são independentes e $\sum_{n=1}^\infty \Prob(A_n) = \infty$, então $\Prob(A_n \text{ i.o.}) = 1$. (Os eventos ruins acontecem infinitas vezes, a convergência falha).
        \end{itemize}
\end{enumerate}

\begin{example}[Exercício 8]
    $X_n \sim U(0, a_n)$. Seja $A_n$ o evento $X_n < 1$. $\Prob(A_n) = \Prob(X_n \le 1) = \frac{1 - 0}{a_n} = 1/a_n$.
    \begin{itemize}
        \item \textbf{a) $a_n = n^2$:} $\Prob(A_n) = 1/n^2$. A soma $\sum \Prob(A_n) = \sum 1/n^2 < \infty$ (Série-p com $p=2 > 1$). Pelo 1º Lema de Borel-Cantelli, $P(A_n \text{ i.o.}) = 0$. "Somente um número finito de $X_n$ toma valores $< 1$".
        \item \textbf{b) $a_n = n$:} $\Prob(A_n) = 1/n$. A soma $\sum \Prob(A_n) = \sum 1/n = \infty$ (Série Harmônica). Pelo 2º Lema de Borel-Cantelli (já que $X_n$ são independentes), $P(A_n \text{ i.o.}) = 1$. "Um número infinito de $X_n$ toma valores $< 1$".
    \end{itemize}
\end{example}

% ---------------------
\subsection{Hierarquia e Relações}
\label{sec:hierarquia}
A relação entre os tipos de convergência é a parte mais importante da teoria.

$$ (X_n \qcto X) \implies (X_n \pto X) \implies (X_n \dto X) $$

\begin{tip}
    **As voltas NÃO são verdadeiras!** A lista de exercícios foi feita para testar isso.
    \begin{itemize}
        \item \textbf{$D \not\implies P$ (Exemplo Q7):} Mostramos que $X_n \dto X \sim Exp(2)$. Mas $X_n$ e $X$ são independentes. $X_n$ não está "chegando perto" de $X$. $\Prob(|X_n - X| > \epsilon)$ não vai para 0, pois $X_n - X$ converge em distribuição para $Y - Z$, onde $Y, Z \sim Exp(2)$ iid, que não é 0.
        \item \textbf{$P \not\implies q.c.$ (Exemplo Q9):} $P(X_n = 1) = 1/n$, $P(X_n = 0) = 1 - 1/n$.
            \begin{itemize}
                \item \textbf{Convergência em P:} $\Prob(|X_n - 0| > \epsilon) = \Prob(X_n = 1) = 1/n$ (para $\epsilon < 1$). Como $\lim 1/n = 0$, temos $X_n \pto 0$.
                \item \textbf{NÃO q.c.:} $\sum \Prob(X_n = 1) = \sum 1/n = \infty$. Pelo 2º Lema de Borel-Cantelli, $X_n=1$ acontece infinitas vezes. A sequência (ex: 1, 0, 1, 0, 0, 1, 0, 0...) nunca se "fixa" em 0. Portanto, $P(\lim X_n = 0) = 0$. $X_n \not\qcto 0$.
            \end{itemize}
    \end{itemize}
\end{tip}

% ---------------------------------------------------
\section{Leis dos Grandes Números (LLN)}
\label{sec:lln}
As Leis dos Grandes Números dizem respeito à convergência da \textbf{média amostral} ($\Xn$) para a \textbf{média populacional} ($\mu$).

\begin{theorem}[Lei Forte dos Grandes Números (SLLN)]
Se $X_1, ..., X_n$ são iid com $\E|X_i| < \infty$ (a média existe), então:
$$ \Xn = \frac{1}{n} \sum_{i=1}^n X_i \qcto \mu \quad \text{onde } \mu = \E(X_1) $$
\end{theorem}

\begin{theorem}[Lei Fraca dos Grandes Números (WLLN)]
Se $X_1, ..., X_n$ são iid com $\E|X_i| < \infty$, então:
$$ \Xn \pto \mu $$
\end{theorem}

\begin{tip}
    A SLLN é o pilar da estatística. Ela garante que, com amostra suficiente, a média amostral \textbf{será} a média real. A SLLN implica a WLLN (pois $q.c. \implies P$).
\end{tip}

\begin{example}[Exercício 17a]
    $Z_i \sim Exp(1/2)$ iid. Qual o limite de $\Xn$?
    \begin{enumerate}
        \item \textbf{Identificar $\mu$:} $\mu = \E(Z_i) = 1/\lambda = 1/(1/2) = 2$.
        \item \textbf{Verificar Condição:} $\E|Z_i| = 2 < \infty$.
        \item \textbf{Aplicar SLLN:} Pela SLLN, $\Xn \qcto 2$.
    \end{enumerate}
    O limite é 2, e o tipo de convergência é (pelo menos) quase certa.
\end{example}

\begin{example}[Exercício 13 - A "pegadinha" do log]
    $X_i \sim U(0,1)$ iid. Ache o limite q.c. da média geométrica $G_n = (\prod_{k=1}^{n} X_k)^{1/n}$.
    \begin{enumerate}
        \item \textbf{Transformar:} Não podemos aplicar a SLLN a um produto. Mas podemos aplicar ao logaritmo:
        $$ \log(G_n) = \log\left( (\prod X_k)^{1/n} \right) = \frac{1}{n} \sum_{k=1}^n \log(X_k) $$
        \item \textbf{Nova V.A.:} Seja $Y_i = \log(X_i)$. $\log(G_n)$ é a média amostral $\bar{Y}_n$.
        \item \textbf{Aplicar SLLN a $\bar{Y}_n$:} $\bar{Y}_n \qcto \E(Y_1) = \E(\log X_1)$.
        $$ \E(\log X_1) = \int_0^1 \log(x) f(x) dx = \int_0^1 \log(x) \cdot 1 dx = [x \log x - x]_0^1 = -1 $$
        (Usando L'Hopital para $\lim_{x \to 0} x \log x = 0$).
        \item \textbf{Concluir:} $\bar{Y}_n = \log(G_n) \qcto -1$.
        \item \textbf{Desfazer o log:} Como $g(x) = e^x$ é contínua, $G_n = e^{\bar{Y}_n} \qcto e^{-1}$.
    \end{enumerate}
\end{example}

% ---------------------------------------------------
\section{Teorema Central do Limite (TCL)}
\label{sec:clt}
O TCL não fala \textit{para onde} a média converge (isso é a LLN), mas sim \textit{como} ela converge. Ele descreve a \textbf{distribuição} do erro em torno da média.

\begin{theorem}[TCL de Lindeberg-Lévy]
Se $X_1, ..., X_n$ são iid com média $\mu$ e variância finita $\sigma^2 > 0$, então:
$$ Z_n = \frac{\Xn - \mu}{\sigma / \sqrt{n}} = \frac{S_n - n\mu}{\sigma \sqrt{n}} \dto N(0, 1) $$
Onde $S_n = \sum X_i$.
\end{theorem}

\begin{tip}
    **TCL vs LLN:**
    \begin{itemize}
        \item \textbf{LLN:} $\Xn \pto \mu$. O limite é uma \textbf{constante}.
        \item \textbf{TCL:} $\sqrt{n}(\Xn - \mu) \dto N(0, \sigma^2)$. O limite é uma \textbf{distribuição}.
    \end{itemize}
    O TCL diz que, não importa a distribuição original (desde que $\sigma^2 < \infty$), a soma (ou média) de muitas V.A.s iid terá uma distribuição aproximadamente Normal.
\end{tip}

\begin{example}[Exercício 20]
    $X_k \sim Bern(0.4)$ iid. $n=100$. $S_{100} = \sum X_k \sim Bin(100, 0.4)$.
    Aproximar $P(S_{100} = 50)$.
    \begin{enumerate}
        \item \textbf{Parâmetros do TCL:} $n=100$.
        $\mu = \E(X_k) = p = 0.4$.
        $\sigma^2 = \Var(X_k) = p(1-p) = 0.4 \cdot 0.6 = 0.24$.
        \item \textbf{Parâmetros da Soma $S_n$:}
        $\E(S_n) = n\mu = 100 \cdot 0.4 = 40$.
        $\Var(S_n) = n\sigma^2 = 100 \cdot 0.24 = 24$.
        \item \textbf{TCL:} $S_{100} \approx N(40, 24)$.
        \item \textbf{Correção de Continuidade (Importante!):}
        $P(S_{100} = 50) \approx P(49.5 < S_{100} < 50.5)$.
        \item \textbf{Padronizar (Z-score):}
        $$ P\left( \frac{49.5 - 40}{\sqrt{24}} < Z < \frac{50.5 - 40}{\sqrt{24}} \right) $$
        $$ P\left( \frac{9.5}{4.899} < Z < \frac{10.5}{4.899} \right) \approx P(1.94 < Z < 2.14) $$
        $$ \Phi(2.14) - \Phi(1.94) \approx 0.9838 - 0.9738 = 0.01 $$
    \end{enumerate}
\end{example}

\begin{tip}[Exercícios 10, 18, 19]
    A lista é inteligente aqui.
    \begin{itemize}
        \item \textbf{Q10:} $X_i \sim N(0,1)$. $S_n \sim N(0, n)$. $S_n/\sqrt{n} \sim N(0, n/(\sqrt{n})^2) = N(0,1)$. Isso é uma distribuição \textbf{exata}, não um limite.
        \item \textbf{Q18:} $X_i \sim N(0,2)$. $S_n \sim N(0, 2n)$. $Y_n = S_n/\sqrt{2n} \sim N(0, 2n/(\sqrt{2n})^2) = N(0,1)$. De novo, \textbf{exata}. A prova por FC (Q18a) vai mostrar $\phi_{Y_n}(t) = e^{-t^2/2}$ para \textit{todo} $n$, não só no limite.
        \item \textbf{Q19:} $E(X_1)=0, E(X_1^2)=2$. Aqui \textbf{não} sabemos se $X_1$ é Normal. \textit{Aqui} sim usamos o TCL. O limite é $N(0,1)$.
    \end{itemize}
    Moral: Não use o TCL se a distribuição da soma já for conhecida e simples (ex: soma de Normais é Normal, soma de Poissons é Poisson).
\end{tip}

% ---------------------------------------------------
\section{O Método Delta}
\label{sec:delta}
O Método Delta é o que usamos quando queremos a distribuição de uma \textit{função} da média amostral, $g(\Xn)$.

\begin{theorem}[Método Delta - 1ª Ordem]
Se $\sqrt{n}(Y_n - \theta) \dto N(0, \tau^2)$ e $g(x)$ é uma função diferenciável em $\theta$ com $g'(\theta) \neq 0$, então:
$$ \sqrt{n}(g(Y_n) - g(\theta)) \dto N(0, \tau^2 [g'(\theta)]^2) $$
\end{theorem}

\begin{tip}
    Na prática, $Y_n = \Xn$, $\theta = \mu$, e $\tau^2 = \sigma^2$. O Método Delta usa uma aproximação de Taylor de 1ª ordem. A nova variância é a variância antiga "esticada" ou "encolhida" pelo quadrado da derivada.
\end{tip}

\subsection{Passo a Passo para resolver}

\begin{enumerate}
    \item \textbf{TCL Básico:} Encontre o limite de $\sqrt{n}(\Xn - \mu) \dto N(0, \sigma^2)$.
    \item \textbf{Identificar $g$ e $\theta$:} Qual a função $g(\cdot)$ e qual o ponto $\theta = \mu$?
    \item \textbf{Derivar:} Calcule a derivada $g'(x)$ e seu valor em $\theta$, $g'(\theta)$.
    \item \textbf{Aplicar a Fórmula:} A distribuição limite é $N(0, \sigma^2 [g'(\theta)]^2)$.
\end{enumerate}

\begin{example}[Exercício 17c]
    Ache o limite de $U_n = \sqrt{n}(\bar{Z}_n^2 - 4)$, onde $Z_i \sim Exp(1/2)$.
    \begin{enumerate}
        \item \textbf{TCL Básico:} De (17a), $\mu = 2$. A variância de $Exp(\lambda)$ é $1/\lambda^2$.
        $\sigma^2 = \Var(Z_i) = 1 / (1/2)^2 = 4$.
        Pelo TCL, $\sqrt{n}(\Xn - 2) \dto N(0, 4)$.
        (Aqui, $\theta = 2$ e $\tau^2 = 4$).
        \item \textbf{Identificar $g$:} A expressão é $\sqrt{n}(g(\Xn) - g(2))$.
        Vemos $g(\Xn) = \Xn^2$ e $g(2) = 2^2 = 4$. A função é $g(x) = x^2$.
        \item \textbf{Derivar:} $g'(x) = 2x$.
        Em $\theta=2$, $g'(2) = 2(2) = 4$.
        \item \textbf{Aplicar Fórmula:} A variância limite é $\tau^2 [g'(\theta)]^2 = 4 \cdot [4]^2 = 64$.
    \end{enumerate}
    Portanto, $\sqrt{n}(\bar{Z}_n^2 - 4) \dto N(0, 64)$.
\end{example}

% ---------------------------------------------------
\section{Estratégia de Resolução: Um Mapa Mental}
\label{sec:estrategia}

\begin{enumerate}
    \item \textbf{Leia o Objeto:} O que está sendo analisado?
    \begin{itemize}
        \item Uma V.A. $X_n$ que muda (ex: $U(0, 1+1/n)$ em \textbf{Q3}, $Exp(2n/n+1)$ em \textbf{Q7})?
        \item Uma função de uma V.A. $g(X_n)$ (ex: $X_n / \log n$ em \textbf{Q1})?
        \item Uma média amostral $\Xn$ (ex: \textbf{Q17a})?
        \item Uma soma $S_n$ (ex: \textbf{Q20})?
        \item Uma V.A. padronizada $Z_n = \sqrt{n}(\Xn - \mu) / \sigma$ (ex: \textbf{TCL})?
        \item Uma função da média $g(\Xn)$ (ex: $\bar{Z}_n^2$ em \textbf{Q17c}, $\log(\Xn)$ em \textbf{Q21})?
    \end{itemize}

    \item \textbf{Leia a Pergunta:} O que se pede?
    \begin{itemize}
        \item \textbf{"Prove que $X_n \to c$"} (para $c$ constante)?
        $\to$ Provavelmente \textbf{Convergência em Probabilidade (Q1, Q2)}.
        $\to$ Use a Definição $P(|X_n - c| > \epsilon) \to 0$.

        \item \textbf{"Encontre o limite em distribuição de $X_n$"}?
        $\to$ Se $X_n$ não for média/soma $\to$ Use o \textbf{Método da FDA (Q3, Q7)}.
        $\to$ Se $X_n$ for uma soma/média padronizada $\to$ Use o \textbf{TCL (Q10, Q18, Q19)}.

        \item \textbf{"Qual o limite de $\Xn$"?} (pergunta sobre o valor)
        $\to$ Use a \textbf{SLLN (Q17a)}. O limite é $\mu$.

        \item \textbf{"Qual o limite em distribuição de $g(\Xn)$"?}
        $\to$ Use o \textbf{Método Delta (Q17c, Q21)}.

        \item \textbf{"Qual o tipo de convergência?"}
        $\to$ Pense na \textbf{Hierarquia (Q7, Q9)}. A SLLN dá $q.c.$, o TCL dá $D$, a WLLN dá $P$.

        \item \textbf{"Aproxime a probabilidade de $S_n$"}?
        $\to$ Use o \textbf{TCL com correção de continuidade (Q20)}.
    \end{itemize}
\end{enumerate}

\end{document}
%%%%%%%%%%%%%%%%%%%%%%%%%%%%%%%%%%%%%%%%%%%%%%%%%%%%%%%%%
% --- FIM DO DOCUMENTO ---
%%%%%%%%%%%%%%%%%%%%%%%%%%%%%%%%%%%%%%%%%%%%%%%%%%%%%%%%%